\chapter{从模形式到自守表示 I}

% ===== PDF p.011 =====
从本讲开始，我们想为 $\GL_2$ 构造与模形式相伴的自守表示。

先回忆上次讨论过的内容：
\begin{itemize}
\item 设 $G$ 是 $\Q$ 上的约化群，$\A=\Af\times\R$ 是阿代尔环。
\item 我们引入空间
\[
[G]=A_GG(\Q)\backslash G(\A),
\]
称之为\textbf{阿代尔商}（adélic quotient），它具有有限的 $G(\A)$-不变体积。
\item 我们陈述过：朗兰兹纲领的目标是按算术数据分解酉 $G(\A)$-表示
\[
L^2([G]).
\]
\item 在例子 $G=\Gm$ 中我们证明了
\[
L^2(\R_{>0}\Q^\times\backslash\A^\times)\cong\widehat{\bigoplus_{\sigma:\Gal(\overline\Q/\Q)\to\C^\times}}\C\chi_\sigma,
\]
作为「所期望的算术数据来自 $\Gal(\overline\Q/\Q)$ 的 Galois 表示」的一个例证。
\item 随后我们讨论了 $[G]$ 的几何：
\begin{itemize}
\item 取 $K_\infty\subseteq G(\R)$ 为某个极大紧连通子群，
\item 取 $K\subseteq G(\Af)$ 为某个紧开子群。
\item 则 $[G]\to[G]/K$ 是 $K$-挠子（特别地，是一个 profinite 覆盖），$[G]/K\to[G]/KK_\infty$ 是 $K_\infty$-挠子，而当 $K$ 充分小时 $[G]/KK_\infty$ 是实流形；更精确地，
\[
[G]/KK_\infty=\bigsqcup_{i=1}^m\Gamma_i\backslash X
\]
是若干算术流形的无交并，其中 $\Gamma_i\subseteq G(\R)$ 是由 $K$ 决定的（离散）同余子群，而
\[
X:=A_G\backslash G(\R)/K_\infty.
\]
\end{itemize}
\item 此外我们定义：$G(\A)$ 的不可约酉表示 $\pi$ 称为\textbf{自守的}（在 $L^2$ 意义下），若它同构于 $L^2([G])$ 的某个子商。
\item 最初的例子（当时只是隐含提及）如下：
\begin{itemize}
\item 若 $G=\Gm$，则 $G(\A)$ 的自守表示（在 $L^2$ 意义下）恰是 $\A^\times$ 的酉特征，且这些特征在 $\Q^\times\cdot\R_{>0}$ 上平凡。
\item 若 $G$ 是一般的群，则平凡表示（即 $[G]$ 上的常函数）是自守的，因为 $[G]$ 有有限体积；并且它落在 $L^2([G])$ 的离散部分 $L^2_{\mathrm{disc}}([G])$ 中。
\end{itemize}
\item 对一般的 $G$，存在正交分解
\[
L^2([G])\cong L^2_{\mathrm{disc}}([G])\oplus L^2_{\mathrm{cont}}([G]),
\]
其中 $L^2_{\mathrm{disc}}([G])$ 是 $G(\A)$ 的所有不可约子表示之和的闭包。
% ===== PDF p.012 =====
\item 最后我们提到，Langlands 用 $L^2_{\mathrm{disc}}([M])$ 描述了 $L^2_{\mathrm{cont}}([G])$，其中 $M$ 遍历抛物子群 $P\subseteq G$（$P\neq G$，且定义在 $\Q$ 上）的 Levi 商。
\end{itemize}

今天我们将更详细地讨论 $G=\GL_2$ 的情形。首先，
\[
A_{\GL_2}\cong\R_{>0}
\]
（作为标量矩阵嵌入 $\GL_2(\R)$）。上次我们断言
\[
\GL_2(\Q)\backslash\GL_2(\A)\cong\GL_2(\Z)\backslash(\GL_2(\widehat\Z)\times\GL_2(\R)),
\]
其中
\[
\widehat\Z:=\varprojlim_m\Z/m\cong\prod_p\Z_p.
\]
下面解释它为什么成立，参见 \texttt{[Tho, Theorem 1]}。

群 $\GL_2(\Af)$ 是诸群 $\GL_2(\Q_p)$ 关于紧开子群 $\GL_2(\Z_p)$ 的限制乘积，即
\[
\GL_2(\Af)=\Bigl\{(A_p)_p\in\prod_p\GL_2(\Q_p)\Bigm| A_p\in\GL_2(\Z_p)\ \text{除有限多个 }p\text{ 外}\Bigr\}
\]
（参见 \texttt{[GH19, Proposition 2.3.1]}）。中国剩余定理蕴涵：$\Q\subseteq\Af$（对角嵌入）是稠密的。此外，由素因子分解，
\[
\Af^\times=\Q^\times\cdot\widehat\Z^\times.
\]

群 $\GL_2(\Z_p)$、$\GL_2(\Q_p)$ 由初等矩阵与对角矩阵生成，这蕴涵
\[
\GL_2(\Af)=\GL_2(\Q)\GL_2(\widehat\Z),
\]
从而
\[
\GL_2(\Q)\backslash\GL_2(\Af)\cong\GL_2(\Q)\cap\GL_2(\widehat\Z)\backslash\GL_2(\widehat\Z)\cong\GL_2(\Z)\backslash\GL_2(\widehat\Z).
\]
由此可推出所想要的\znote{此处关键的一步是 $\SL_2$ 的强逼近定理：$\SL_2(\Q)\SL_2(\widehat\Z)=\SL_2(\Af)$，即 $\SL_2(\Q)$ 在 $\SL_2(\Af)$ 中是稠密的；再配合 $\Af^\times=\Q^\times\widehat\Z^\times$ 处理行列式，就得到上面的双陪集分解。原讲义把这一步归为「初等矩阵与对角矩阵生成」，细节可参见 \texttt{[Tho, Theorem 1]}。[来源：原文第 12 页与 \texttt{[Tho, Theorem 1]}]}
\[
\GL_2(\Q)\backslash\GL_2(\A)\cong\GL_2(\Z)\backslash(\GL_2(\widehat\Z)\times\GL_2(\R)).
\]
上述证明其实对一切 $n\ge1$ 都成立，并表明
\[
\GL_n(\Q)\backslash\GL_n(\A)\cong\GL_n(\Z)\backslash(\GL_n(\widehat\Z)\times\GL_n(\R)).
\]

对 $m\ge1$ 定义
\[
K(m)=\ker(\GL_2(\widehat\Z)\to\GL_2(\Z/m)).
\]
这些群在所有紧开子群 $\GL_2(\Af)$ 中共尾。我们有
\[
\GL_2(\Q)\backslash\GL_2(\Af)/K(m)\cong\GL_2(\Z)\backslash\GL_2(\widehat\Z)/K(m),
\]
且 $\GL_2(\widehat\Z)/K(m)\cong\GL_2(\Z/m)$。群 $\SL_2(\Z/m)$ 甚至由初等矩阵生成\znoteplain{原文脚注：对某个环 $R$ 中的任意可逆元 $a$，利用欧几里得算法与如下「神奇恒等式」：
\[
\begin{pmatrix}a&0\\0&a^{-1}\end{pmatrix}
=\begin{pmatrix}1&a\\0&1\end{pmatrix}
\begin{pmatrix}1&0\\-a^{-1}&1\end{pmatrix}
\begin{pmatrix}1&a\\0&1\end{pmatrix}
\begin{pmatrix}0&-1\\1&0\end{pmatrix}.
\]}。这蕴涵态射 $\SL_2(\Z)\to\SL_2(\Z/m)$ 是满射。于是我们得到
\[
\GL_2(\Z)\backslash\GL_2(\widehat\Z)/K(m)\cong\{\pm1\}\backslash(\Z/m)^\times,
\]
% ===== PDF p.013 =====
同构由行列式给出。现在定义
\[
K_\infty:=\SO_2(\R)\subseteq\GL_2(\R)
\]
以及
\[
X:=A_{\GL_2}\backslash\GL_2(\R)/K_\infty.
\]
则 $X\cong\HH^\pm$（上半平面与下半平面之并），同构由态射
\[
A_{\GL_2}\backslash\GL_2(\R)/K_\infty\xrightarrow{\ \sim\ }\HH^\pm,\qquad
g=\begin{pmatrix}a&b\\c&d\end{pmatrix}\mapsto g(i):=\frac{a\cdot i+b}{c\cdot i+d}
\]
给出。于是最终得到
\[
\GL_2(\Q)\backslash(\GL_2(\Af)/K(m)\times X)\cong\bigsqcup_{\{\pm1\}\backslash(\Z/m)^\times}\Gamma(m)\backslash\HH^\pm,
\]
其中
\[
\Gamma(m):=K(m)\cap\GL_2(\Z)=\ker(\GL_2(\Z)\to\GL_2(\Z/m)).
\]

$\GL_2(\Q)\backslash(\GL_2(\Af)/K(m)\times X)$ 不连通这件事也许显得烦人，但后面会表明这其实是优点。注意：若把 $K\subseteq\GL_2(\Af)$ 换成某个使
\[
\det:K\to\widehat\Z^\times
\]
为满射的紧开子群 $K$，例如取
\[
K_1(m):=\Bigl\{\begin{pmatrix}a&b\\c&d\end{pmatrix}\in\GL_2(\Af)\Bigm| c\equiv0,\ d\equiv1\bmod m\Bigr\},
\]
则 $\GL_2(\Q)\backslash(\GL_2(\Af)/K\times X)$ 连通。

$G=\GL_2$ 时 $[G]$ 体积的有限性，现在可以用直接论证来检验。$K_\infty K(m)$ 的紧性蕴涵：只需证明空间
\[
[\GL_2]/K_\infty K(m)
\]
有有限体积。把这个空间写成无交并，甚至可以归约到「$\Gamma(m)\backslash\HH^\pm$ 对 $m\ge1$ 有有限体积」这一陈述。至多差 $\R_{>0}$ 中的一个标量，$\HH^\pm$ 上 $\GL_2(\R)$-不变测度由体积形式
\[
\frac{1}{y^2}\,\mathrm{d}x\wedge\mathrm{d}y
\]
给出。确实，我们须证明对 $g\in\GL_2(\R)$ 有
\[
g^*\Bigl(\frac{1}{y^2}\,\mathrm{d}x\wedge\mathrm{d}y\Bigr)=\frac{1}{y^2}\,\mathrm{d}x\wedge\mathrm{d}y.
\]
现在把 $\mathrm{d}x\wedge\mathrm{d}y$ 写成 $\frac{i}{2}\,\mathrm{d}z\wedge\mathrm{d}\bar z$，其中 $z=x+iy$。容易验证上述陈述对
\[
g=\begin{pmatrix}1&a\\0&1\end{pmatrix},\qquad a\in\R,
\]
对 $A_{\GL_2}$，对
\[
g=\begin{pmatrix}a&0\\0&1\end{pmatrix},\qquad a\in\R^\times,
\]
最后对
\[
g=\begin{pmatrix}0&1\\1&0\end{pmatrix},
\]
即对 $z\mapsto\frac{1}{z}=\frac{\bar z}{|z|^2}$ 成立。由前面提到的神奇恒等式，可以推出对下三角矩阵的不变性，从而对 $\GL_2(\R)$ 全体成立。

% ===== PDF p.014 =====
现在证明 $\Gamma(m)\backslash\HH^\pm$ 有有限体积。$\SL_2(\Z)$ 在 $\HH$ 上的一个已知基本域由集合
\[
D:=\{z\in\HH\mid |z|>1,\ |\operatorname{Re}(z)|<1/2\}
\]
给出，而 $\GL_2(\Z)\backslash\HH^\pm$ 体积的有限性由
\[
\int_D\frac{1}{y^2}\,\mathrm{d}x\wedge\mathrm{d}y<\infty
\]
推出。最后，由于 $\Gamma(m)\subseteq\GL_2(\Z)$ 是有限指数的，$\Gamma(m)\backslash\HH^\pm$ 有有限体积，如所愿。

朗兰兹纲领尤其包含模形式与 $\Gal(\overline\Q/\Q)$ 的二维（$\ell$-进）Galois 表示之间的紧密联系；讨论这一联系将是本课程的一大主题。首先回忆模形式的经典定义，参见 \texttt{[DS05]}。为此，设 $\Gamma\subseteq\SL_2(\Z)$ 为同余子群（后面假设它充分小），并固定 $k\in\Z$。

\noindent\textbf{定义 3.1}\phantomsection\label{def:3.1}\quad 函数 $f:\HH^\pm\to\C$ 称为 $\Gamma$ 的\textbf{权 $k$ 模形式}，若
\begin{itemize}
\item 对 $\gamma=\begin{pmatrix}a&b\\c&d\end{pmatrix}\in\Gamma$ 有 $f(\gamma\cdot z)=(cz+d)^kf(z)$；
\item $f$ 全纯；
\item $f$ 在 $\Gamma$ 的尖点处全纯（稍后解释）。
\end{itemize}
令 $M_k(\Gamma)$ 为 $\Gamma$ 的权 $k$ 模形式构成的 $\C$-向量空间。只满足前两个条件、而不必在尖点处全纯的函数，我们也称为\textbf{弱模形式}。

回忆：上半平面 $\HH$ 自然 $\SL_2(\Q)$-等变地嵌入集合
\[
\HH^*:=\HH\cup\Pone(\Q),
\]
后者带 Satake 拓扑。同余子群 $\Gamma\subseteq\SL_2(\Z)$ 在 $\Pone(\Q)$ 上的轨道称为 $\Gamma$ 的\textbf{尖点}，商
\[
\Gamma\backslash\HH^*
\]
是 $\Gamma\backslash\HH$ 的自然紧化。把 $\Pone(\Q)$ 的另一份加到下半平面 $\HH^-$ 上，对 $\Gamma\subseteq\GL_2(\Z)$ 就得到类似的对象与概念。

「在尖点处全纯」这一条件现在意味着：对每个尖点 $\sigma$，$\Gamma$ 的弱模形式 $f$ 都有 Fourier 展开，即把该尖点平移到 $\infty\in\Pone(\Q)$ 之后，存在等式
\[
f(z)=\sum_{n\in\Z}a_nq^{n/h},
\]
其中 $q=e^{2\pi iz}$，$h$ 是所谓 $f$ 的\textbf{最小周期}，$a_n\in\C$，$n\in\Z$。若 $n<0$ 时 $a_n=0$，则称弱模形式 $f$ 在该尖点处全纯；若它在每个尖点处都全纯，则称它在尖点处全纯。若进一步对每个尖点都有 $a_0=0$，则称 $f$ 为\textbf{尖点形式}。记
\[
S_k(\Gamma)\subseteq M_k(\Gamma)
\]
为 $\Gamma$ 的权 $k$ 尖点形式子空间。

下一个目标是把 $\Gamma\subseteq\GL_2(\Z)$ 的权 $k\in\Z$ 的（弱）模形式典范地描述为某个全纯线丛
% ===== PDF p.015 =====
$\omega^{\otimes k}$ 在 $\Gamma\backslash\HH^\pm$ 上（至少当 $\Gamma$ 充分小时）的截面。

做这件事之前，我们先讨论一些有用的「纤维丛的下降」构造，它们会被频繁使用。

考虑一个集合 $Z$，以及一个自由右作用在 $Z$ 上的群 $H$。则投影
\[
Z\to Z/H
\]
配得上称为 $H$-\textbf{挠子}。设 $F$ 是带 $H$（左）作用的集合。$Z$ 关于 $H$ 与 $F$ 的\textbf{收缩积}（contracted product）
\[
Z\times^HF
\]
定义为 $Z\times F$ 对 $H$ 的作用 $h\cdot(z,f)\mapsto(zh^{-1},hf)$ 的商。经由自然投影 $(z,f)\mapsto z$，我们得到 $Z/H$ 上的「纤维丛」
\[
Z\times^HF\to Z/H,
\]
其纤维（非典范地）同构于 $F$。确实，对类 $c\in Z/H$ 与满足 $z\mapsto c$ 的 $z\in Z$ 的选取，得到同构
\[
F\cong(Z\times^HF)\times_{Z/H}c,\qquad f\mapsto(z,f).
\]
注意 $-\times^H-$ 表示收缩积，而 $-\times_{Z/H}-$ 表示纤维积（这是两个互不相干的构造）。特别地，存在典范同构
\[
Z\times F\cong(Z\times^HF)\times_{Z/H}Z,\qquad(z,f)\mapsto((z,f),z),
\]
即纤维丛 $Z\times^HF$ 拉回到 $Z$ 之后典范地平凡化。\znoteplain{原文此处说拉回到 $Z/H$；但按上式应为拉回到 $Z$（$Z/H$ 是底空间）。}

我们可以方便地描述纤维丛 $Z\times^HF\to Z/H$ 的截面。

\noindent\textbf{引理 3.2}\phantomsection\label{lem:3.2}\quad 截面 $s:Z/H\to Z\times^HF$（即 $Z/H$ 上的态射）的集合，典范地等同于满足
\[
\phi(zh)=h^{-1}\phi(z)
\]
（对一切 $z\in Z$，$h\in H$）的函数 $\phi:Z\to F$ 的集合。

\noindent\textbf{证明}．给定 $s:Z/H\to Z\times^HF$ 与 $z\in Z$，由 $H$ 在 $Z$ 上作用的自由性，存在唯一元素 $f\in F$ 使
\[
s(z)=(z,f).
\]
令 $\phi(z):=f$。则
\[
(z,\phi(z))=s(z)=s(zh)=(zh,\phi(zh))=(z,h\phi(zh)),
\]
这蕴涵
\[
\phi(z)=h\phi(zh)
\]
如所愿。反之，给定满足 $\phi(zh)=h^{-1}\phi(z)$（对一切 $z\in Z$，$h\in H$）的函数 $\phi:Z\to F$，函数
\[
Z\to Z\times F,\qquad z\mapsto(z,\phi(z))
\]
按定义对 $H$ 的作用不变，从而下降为截面 $s:Z/H\to Z\times^HF$。两个构造互为逆。\hfill$\square$

% ===== PDF p.016 =====
反之，假设给定态射 $Y\to Z/H$，以及 $Z$ 上的同构
\[
g:Y\times_{Z/H}Z\xrightarrow{\ \sim\ }Z\times F
\]
（对某个集合 $F$）。取 $z\in Z$，其像为 $c\in Z/H$。则对任意 $h\in H$ 得到同构
\[
\sigma_h:F\xrightarrow{\ g^{-1}|_{z\times F}\ }Y_c\times_cz\xrightarrow{\ \cdot h\ }Y_c\times_czh\xrightarrow{\ g|_{g^{-1}(zh\times F)}\ }F,
\]
它们定义了 $H$ 在 $F$ 上的作用。此外可以验证，通过这个作用
\[
Y\cong Z\times^HF.
\]
用更花哨的话说：带 $Z$ 上「平凡化」的「纤维丛」$Y\to Z/H$，等价于带 $H$ 作用的集合。

我们再提一个变体：当 $Z=G$ 是群、$H\subseteq G$ 是子群时。$G/H$ 上的 $G$-\textbf{等变集合} $Y$ 是 $G/H$ 上的空间 $\pi:Y\to G/H$，连同 $G$ 的左作用，使 $\pi$ 为 $G$-等变。此时有典范平凡化
\[
G\times^{1\cdot H}Y\cong Y\times_{G/H}G,\qquad(g,y)\mapsto(gy,g),
\]
由此可推出：$G/H$ 上 $G$-等变集合的范畴恰好等价于带 $H$（左）作用的集合的范畴，因为对 $H$-集合 $F$，收缩积
\[
G\times^HF
\]
自然地是 $G$-等变的（让 $G$ 通过左乘作用在左因子上）。注意我们要求 $H$ 右作用在 $Z$（相应地 $G$）上，但当然也可以类似地考虑 $H\backslash Z$ 上的「纤维丛」。

我们把构造放在集合范畴中，但显然若 $Z$、$H$、$F$ 是拓扑空间（实流形、复流形……）且 $H$ 的作用连续（光滑、全纯……），我们可以在拓扑空间（实流形、复流形……）等范畴中作同样的构造。

现在开始把模形式描述为模曲线\idx{模曲线}上的线丛的截面。

首先注意，线丛
\[
\OO_{\Pone_\C}(k)
\]
（$k\in\Z$）自然地是 $\GL_2(\C)$-等变的。事实上，取张量幂后只需考虑 $k=1$。但 $\OO_{\Pone_\C}(1)$ 相伴于自然的 $\Gm$-挠子
\[
\mathbb{A}^2_\C\setminus\{0\}\to\Pone_\C,\qquad(z_1,z_2)\mapsto[z_1:z_2],
\]
后者自然地 $\GL_2(\C)$-等变。其次注意，嵌入
\[
\HH^\pm\subseteq\Pone_\C,\qquad z\mapsto[z:1]
\]
是 $\GL_2(\R)$-等变的，于是对 $k\in\Z$ 得到 $\GL_2(\R)$-等变线丛
\[
\OO_{\HH^\pm}(k):=\OO_{\Pone_\C}(k)\times_{\Pone_\C}\HH^\pm,
\]
它位于 $\HH^\pm$ 上。

% ===== PDF p.017 =====
$i\in\HH^\pm$ 在 $\GL_2(\R)$ 中的稳定化子等同于
\[
\C^\times\cong\R_{>0}\SO_2(\R)\subseteq\GL_2(\R).
\]
利用
\[
\HH^\pm\cong\GL_2(\R)/\C^\times,
\]
在先前提到的纤维丛构造下，$\OO_{\HH^\pm}(k)$ 对应于（全纯）表示
\[
\chi_k:\C^\times\to\GL_1(\C)=\C^\times,\qquad z\mapsto z^k.
\]
注意引理 3.2 蕴涵：$\OO_{\HH^\pm}(k)$ 在 $\HH^\pm$ 上的光滑截面等同于满足
\[
\phi(gz)=z^{-k}\phi(g),\qquad g\in\GL_2(\R),\ z\in\C^\times
\]
的光滑函数 $\phi:\GL_2(\R)\to\C$。这一点后面会用到。从现在起我们假设 $\Gamma$ 充分小，即它自由作用在 $\HH^\pm$ 上。对 $\Gamma$ 在 $\OO_{\HH^\pm}(k)$ 上的作用取商，得到 $\Gamma\backslash\HH^\pm$ 上的全纯线丛
\[
\omega^{\otimes k}:=\Gamma\backslash\OO_{\HH^\pm}(k).
\]
为把 $\omega^{\otimes k}$ 的截面与 $\Gamma$ 的权 $k$（弱）模形式联系起来，我们须把 $\GL_2(\R)$ 在 $\OO_{\HH^\pm}(k)$ 上的作用写得更显式。对 $k\in\Z$ 定义 $\HH^\pm$ 上的全纯 $\GL_2(\R)$-等变线丛
\[
L_k:=\HH^\pm\times\C,
\]
$\GL_2(\R)$ 通过
\[
g\cdot(z,\lambda):=\Bigl(\frac{az+b}{cz+d},(cz+d)^k\lambda\Bigr)
\]
左作用，其中 $g=\begin{pmatrix}a&b\\c&d\end{pmatrix}\in\GL_2(\R)$。为记号方便，引入函数（一个「自守因子」）\znote{自守因子 $j(g,z)=cz+d$ 有简单的几何意义：$\operatorname{Im}(gz)=\det(g)|j(g,z)|^{-2}\operatorname{Im}(z)$，且它满足余圈关系 $j(gh,z)=j(g,hz)j(h,z)$。权 $k$ 的变换律正是用 $j$ 的 $k$ 次幂写出的。[来源：原文第 17、25 页]}
\[
j:\GL_2(\R)\times\HH^\pm\to\C^\times,\qquad
\Bigl(g=\begin{pmatrix}a&b\\c&d\end{pmatrix},z\Bigr)\mapsto j(g,z):=cz+d.
\]
留作练习：验证映射
\[
L_1^*\to\OO_{\HH^\pm}(1)^*=\HH^\pm\times_{\Pone_\C}\mathbb{A}^2_\C\setminus\{0\},\qquad
(z,\lambda)\mapsto(z,(\lambda z,\lambda))
\]
是 $\HH^\pm$ 上 $\GL_2(\R)$-等变线丛的全纯同构。这里（以及下文）$(-)^*$ 表示线丛中零截面的补集。取张量幂得到
\[
L_k\cong\OO_{\HH^\pm}(k)
\]
% ===== PDF p.018 =====
对一切 $k\in\Z$ 成立，作为 $\GL_2(\R)$-等变全纯线丛\znoteplain{原文脚注：$L_k$ 与 $\OO_{\HH^\pm}(k)$ 作为 $\GL_2(\R)$-等变线丛抽象同构，也可由观察到两者都对应于 $\C^\times$-表示 $z\mapsto z^k$ 推出。这一点已对 $\OO_{\HH^\pm}$ 指出，其成立是因为在嵌入 $\C^\times\hookrightarrow\GL_2(\R)$ 下，元素 $z\in\C^\times$ 映到矩阵 $g=\begin{pmatrix}\operatorname{Re}(z)&-\operatorname{Im}(z)\\\operatorname{Im}(z)&\operatorname{Re}(z)\end{pmatrix}$，且 $j(g,i)=z$。}。

与引理 3.2 类似\znoteplain{原文脚注：这里的情形与引理 3.2 略有不同，即这些 $\Gamma$-等变向量丛并不相伴于 $\Gamma$ 的某个表示。}，我们得到：全纯截面
\[
\HH^\pm\to L_k,\qquad z\mapsto(z,f(z))
\]
是 $\Gamma$-等变的，当且仅当 $f(z)$ 满足
\[
f(\gamma\cdot z)=j(\gamma,z)^kf(z)
\]
（$\gamma\in\Gamma$），即
\[
f\Bigl(\frac{az+b}{cz+d}\Bigr)=(cz+d)^kf(z),
\]
只要 $\gamma=\begin{pmatrix}a&b\\c&d\end{pmatrix}\in\Gamma$。换言之，我们证明了
\[
H^0(\Gamma\backslash\HH^\pm,\omega^{\otimes k})
\]
等同于 $\Gamma$ 的弱模形式空间。特别地，有自然嵌入
\[
M_k(\Gamma)\subseteq H^0(\Gamma\backslash\HH^\pm,\omega^{\otimes k}),
\]
其像由「在尖点处全纯」这一条件定义。下面分析在尖点处全纯对 $\omega^{\otimes k}$ 的截面意味着什么。

我们断言 $\omega^{\otimes k}$ 典范地延拓到 $\Gamma\backslash\HH^\pm$ 的（Satake）紧化 $\overline{\Gamma\backslash\HH^\pm}$。尖点 $\infty\in\Pone(\Q)$ 在 $\SL_2(\Z)$ 中的稳定化子 $\Gamma_\infty$ 是元素
\[
\begin{pmatrix}1&a\\0&1\end{pmatrix},\qquad a\in\Z
\]
构成的子群。因此只需看 $\omega$ 沿嵌入
\[
\Gamma_\infty\backslash(\HH\cup\{\infty\})
\]
的（典范）延拓。记
\[
L_k^+:=L_k|_{\HH}.
\]
则 $\Gamma_\infty$ 在 $L_k^+=\HH\times\C$ 上的作用只通过第一个因子（因为对一切 $z\in\HH$ 与任意 $\gamma\in\Gamma_\infty$ 有 $j(\gamma,z)=1$）。注意经由指数映射
\[
\Gamma_\infty\backslash\HH\cong\mathbb{D}^\times:=\{q\in\C^\times\mid 0<|q|<1\}
\]
（讨论 Fourier 展开时已隐含使用了这一点）。于是典范地
\[
\Gamma_\infty\backslash L_k^+\cong\mathbb{D}^\times\times\C,
\]
它显然延拓到 $\mathbb{D}:=\{q\in\C\mid|q|<1\}$，即延拓为 $\mathbb{D}\times\C$。这便证明了 $\omega^{\otimes k}$ 典范地延拓到 $\Gamma\backslash\HH^\pm$ 的上述（Satake）紧化 $\overline{\Gamma\backslash\HH^\pm}$。我们把这个延拓仍记作 $\omega^{\otimes k}$。全纯性条件
% ===== PDF p.019 =====
在尖点处全纯，现在翻译为：截面 $f\in H^0(\Gamma\backslash\HH^\pm,\omega^{\otimes k})$ 延拓为 $\overline{\Gamma\backslash\HH^\pm}$ 上 $\omega^{\otimes k}$ 的截面\znote{「在尖点处全纯」等价于 $q$-展开 $f=\sum_{n\in\Z}a_nq^{n/h}$ 中不出现负幂次；而「在尖点处全纯且常数项为零」则进一步要求 $a_0=0$。几何上，前者恰是截面沿紧化边界全纯延拓，后者是截面在边界除子 $D$ 上消没。[来源：原文第 14、19 页]}。

再多用一点理论：Riemann 曲面 $\Gamma\backslash\HH^\pm$ 是 $\C$ 上的代数曲线，其典范紧化\idx{典范紧化}的解析化是 $\overline{\Gamma\backslash\HH^\pm}$。由于 $M_k(\Gamma)$ 是 $H^0(\Gamma\backslash\HH^\pm,\omega^{\otimes k})$ 在 $H^0(\overline{\Gamma\backslash\HH^\pm},\omega^{\otimes k})$ 中的像，由 GAGA 可推出 $M_k(\Gamma)$ 在 $\C$ 上有限维。此外可以证明 $\omega$ 在 $\Gamma\backslash\HH^\pm$ 上丰沛，这蕴涵 $k<0$ 时 $M_k(\Gamma)=0$，而 $k\gg0$ 时 $M_k(\Gamma)$ 变得「大」。事实上，用 Riemann--Roch 定理可以计算 $M_k(\Gamma)$ 的维数。最后注意
\[
S_k(\Gamma)\cong H^0(\Gamma\backslash\HH^\pm,\omega^{\otimes k}(-D)),
\]
其中 $D=\overline{\Gamma\backslash\HH^\pm}\setminus\Gamma\backslash\HH^\pm$ 是无穷远处的（既约）除子。更多细节参见 \texttt{[DS05, Chapter 3, Chapter 7]}。

现在解释如何从 $\omega^{\otimes k}$ 的截面过渡到 $\GL_2(\A)$（或 $[\GL_2]$）上的函数。为此，我们把 $\Gamma\subseteq\GL_2(\Z)$ 遍历所有同余子群时的模形式，解释为 $K\subseteq\GL_2(\Af)$ 遍历所有紧开子群时的「模形式」。

设 $K\subseteq\GL_2(\Af)$ 为紧开子群，并写
\[
\GL_2(\Q)\backslash(\GL_2(\Af)/K\times\HH^\pm)\cong\bigsqcup_{i=1}^m\Gamma_i\backslash\HH^\pm,
\]
其中 $\Gamma_i\subseteq\GL_2(\Z)$ 是某些同余子群。则由 $\HH^\pm$ 拉回，$\OO_{\HH^\pm}(k)$ 定义了
\[
\GL_2(\Q)\backslash(\GL_2(\Af)/K\times\HH^\pm)
\]
上的（复）线丛 $\omega^{\otimes k}$（我们隐含假设 $K$ 充分小，并用了 $\OO_{\HH^\pm}(k)$ 的 $\GL_2(\Q)$-等变性）。

$\Gamma_i$（$i=1,\dots,m$）的模形式空间 $M_k(\Gamma_i)$（甚至弱模形式）嵌入到全纯截面空间
\[
H^0(\GL_2(\Q)\backslash(\GL_2(\Af)/K\times\HH^\pm),\omega^{\otimes k}).
\]
令
\[
M_k(K):=\bigoplus_{i=1}^mM_k(\Gamma_i).
\]
则 $M_k(K)\subseteq H^0(\GL_2(\Q)\backslash(\GL_2(\Af)/K\times\HH^\pm),\omega^{\otimes k})$ 同样由「在所有尖点处全纯」这一条件定义。我们甚至可以去掉 $K$，用 $\omega^{\otimes k}$ 的拉回定义
\[
\GL_2(\Q)\backslash(\GL_2(\Af)\times\HH^\pm)
\]
上的 $G(\Af)$-等变线丛，仍记作 $\omega^{\otimes k}$。注意这个空间不再是复流形，而只是复流形沿有限覆盖映射的逆向极限。但仍可通过
\[
H^0(\GL_2(\Q)\backslash(\GL_2(\Af)\times\HH^\pm),\omega^{\otimes k})
:=\varinjlim_KH^0(\GL_2(\Q)\backslash(\GL_2(\Af)/K\times\HH^\pm),\omega^{\otimes k})
\]
定义它上面 $\omega^{\otimes k}$ 的「全纯截面」的类似物。

% ===== PDF p.020 =====
此外，我们得到（大空间）
\[
M_k:=\varinjlim_KM_k(K),
\]
即权 $k$ 的模形式。\znoteplain{原文写作 “modular forms of level $k$”，其中 “level” 应为 “weight”（级与权是两个概念）。}后面我们会解释：对（充分小的）$K\subseteq\GL_2(\Af)$，$M_k(K)$ 是 $M_k$ 上某个典范 $K$-作用的 $K$-不变量。

由围绕引理 3.2 的纤维丛讨论我们知道：$\OO_{\HH^\pm}(k)$ 沿
\[
\GL_2(\R)\to\HH^\pm,\qquad g\mapsto g\cdot i
\]
的拉回典范平凡（作为 $\GL_2(\R)$-等变丛），且 $\OO_{\HH^\pm}(k)$ 在 $\HH^\pm$ 上的光滑截面典范等同于满足
\[
\phi(gz)=z^{-k}\phi(g),\qquad g\in\GL_2(\R),\ z\in\C^\times
\]
的光滑函数 $\phi:\GL_2(\R)\to\C$。具体地，把光滑截面
\[
\HH^\pm\to L_k=\HH^\pm\times\C,\qquad z\mapsto(z,f(z))
\]
对应到光滑函数
\[
\phi_f:\GL_2(\R)\to\C,\qquad g\mapsto j(g,i)^{-k}f(gi).
\]
注意 $L_k$ 在 $\HH^\pm$ 上的光滑截面就是 $\HH^\pm$ 上的光滑函数（但这个等同不是 $\GL_2(\R)$-等变的！）。

若把光滑性概念略加推广，我们可以把 $\HH^\pm$ 换成 $\GL_2(\Af)\times\HH^\pm$ 做类似的事。

\noindent\textbf{定义 3.3}\phantomsection\label{def:3.3}\quad 设 $Y_i$（$i\in I$）是实流形的共滤逆向系统，使每个过渡映射 $Y_i\to Y_j$ 都是有限覆盖。令 $Y:=\varprojlim_{i\in I}Y_i$，$\pi_i:Y\to Y_i$ 为典范投影。则称函数
\[
\phi:Y\to\C
\]
为\textbf{光滑的}，若存在 $i\in I$ 与光滑函数 $\phi_i:Y_i\to\C$ 使 $\phi=\phi_i\circ\pi_i$。

注意按同样的模式，若每个 $Y_i$ 还是复流形，也可定义 $Y$ 上的全纯函数（或线丛截面）。

当然，我们想把这个术语用到 $\GL_2(\Af)\times\HH^\pm$（或 $\GL_2(\A)=\GL_2(\Af)\times\GL_2(\R)$）上，它是复流形 $\GL_2(\Af)/K\times\HH^\pm$（相应地实流形 $\GL_2(\Af)/K\times\GL_2(\R)$）的逆向极限，其中 $K$ 遍历 $\GL_2(\Af)$ 的紧开子群的共滤系统。我们得到：光滑截面
\[
\GL_2(\Af)\times\HH^\pm\to\GL_2(\Af)\times L_k,\qquad(g,z)\mapsto(g,z,f(g,z))
\]
等同于满足
\[
\phi(g,g_\infty z)=z^{-k}\phi(g,g_\infty)
\]
% ===== PDF p.021 =====
（对 $(g,g_\infty)\in\GL_2(\A)$ 与 $z\in\C^\times\subseteq\GL_2(\R)$）的光滑函数
\[
\phi:\GL_2(\Af)\times\GL_2(\R)\to\C.
\]
具体地，对 $f$ 令
\[
\phi_f(g,g_\infty):=j(g_\infty,i)^{-k}f(g,g_\infty i).
\]
注意截面
\[
\GL_2(\Af)\times\HH^\pm\to\GL_2(\Af)\times L_k,\qquad(g,z)\mapsto(g,z,f(g,z))
\]
在 $\GL_2(\Q)$ 的左作用下不变，当且仅当 $f$ 满足模性条件
\[
f(\gamma g,\gamma z)=j(\gamma,z)^kf(g,z).
\]
利用等式
\[
j(\gamma g_\infty,i)=j(\gamma,g_\infty i)j(g_\infty,i),
\]
我们得到：$f$ 定义 $L_k$ 的 $\GL_2(\Q)$-等变截面，当且仅当函数 $\phi_f$ 是 $\GL_2(\Q)$-不变的，即
\[
\phi_f(\gamma g,\gamma g_\infty)=\phi_f(g,g_\infty)
\]
对 $(g,g_\infty)\in\GL_2(\A)$ 成立。

限制到全纯截面，我们于是构造了自然态射
\[
H^0(\GL_2(\Q)\backslash(\GL_2(\Af)\times\HH^\pm),\omega^{\otimes k})\to C^\infty(\GL_2(\Q)\backslash\GL_2(\A))
\]
使 $f(g,z)\mapsto\phi_f(g,g_\infty):=j(g_\infty,i)^{-k}f(g,g_\infty i)$；这里我们把 $\GL_2(\Q)\backslash(\GL_2(\Af)\times\HH^\pm)$ 上 $\omega^{\otimes k}$ 的截面看作满足上述模性条件的函数
\[
\GL_2(\Af)\times\HH^\pm\to\C,\qquad(g,z)\mapsto f(g,z).
\]

回忆
\[
M_k=\varinjlim_KM_k(K),
\]
其中 $M_k(K)$ 是各种同余子群 $\Gamma_i\subseteq\GL_2(\Z)$ 的模形式空间之和
\[
\bigoplus_{i=1}^mM_k(\Gamma_i),
\]
而 $M_k$ 在 $H^0(\GL_2(\Q)\backslash(\GL_2(\Af)\times\HH^\pm),\omega^{\otimes k})$ 中的像由「在尖点处全纯」这一条件定义。

这样，我们就把模形式与 $\GL_2(\Q)\backslash\GL_2(\A)$ 上的光滑函数联系了起来。注意
\[
[\GL_2]=\GL_2(\Q)A_{\GL_2}\backslash\GL_2(\A),
\]
但我们的函数 $\phi_f$ 不在 $A_{\GL_2}$ 下不变，即
\[
\phi(g,g_\infty r)=r^{-k}\phi(g,g_\infty),\qquad r\in A_{\GL_2}\cong\R_{>0}.
\]
不过这不难修正。令
\[
|\cdot|_{\mathrm{ad\acute{e}lique}}:\A^\times\to\R_{>0},\qquad((x_p)_p,x_\infty)\mapsto\prod_p|x_p|_p|x_\infty|
\]
为\textbf{adélic 绝对值}（它在 $\Q^\times$ 上平凡）。则对 $f\in M_k$，函数
\[
\widetilde\phi_f(g,g_\infty):=|\det(g,g_\infty)|^{k/2}_{\mathrm{ad\acute{e}lique}}\phi_f(g,g_\infty)
\]
在 $\GL_2(\Q)A_{\GL_2}$ 下不变，即定义了光滑函数 $[\GL_2]\to\C$。

% ===== PDF p.022 =====
我们必须澄清模形式何时给出 $L^2([\GL_2])$ 中的函数。为此取截面 $f\in H^0(\GL_2(\Q)\backslash(\GL_2(\Af)\times\HH^\pm),\omega^{\otimes k})$。问题是：
\[
\widetilde\phi_f\in C^\infty([\GL_2])
\]
何时属于 $L^2([\GL_2])$？取紧开子群 $K\subseteq\GL_2(\Af)$ 使
\[
f\in H^0(\GL_2(\Q)\backslash(\GL_2(\Af)/K\times\HH^\pm),\omega^{\otimes k}).
\]
函数 $\widetilde\phi_f$ 是 $L^2$ 的，当且仅当光滑函数 $\widetilde\phi_f\in C^\infty([\GL_2]/K)\subseteq C^\infty([\GL_2])$ 是 $L^2$ 的（因为 $K$ 紧）。这把问题归约到：对同余子群 $\Gamma\subseteq\GL_2(\Z)$ 与 $f\in H^0(\Gamma\backslash\HH^\pm,\omega^{\otimes k})$，函数
\[
\Gamma A_{\GL_2}\backslash\GL_2(\R)\to\C,\qquad
g\mapsto\widetilde\phi_f(g)=|\det(g)|^{k/2}(ci+d)^{-k}f(g(i))
\]
何时平方可积，即属于 $L^2$？我们断言：
\[
|\widetilde\phi_f(g)|^2=|f(g(i))|^2|\operatorname{Im}(g(i))|^k
\]
对任意 $g\in\GL_2(\R)$ 成立，等价地，
\[
|\det(g)|^{1/2}|ci+d|^{-1}=|\operatorname{Im}(g(i))|^{1/2}
\]
对 $g=\begin{pmatrix}a&b\\c&d\end{pmatrix}\in\GL_2(\R)$ 成立。\znoteplain{原文此式写作 $|\det(g)|^{-1/2}|ci+d|=|\operatorname{Im}(g(i))|^{1/2}$，与上一式互为倒数（取 $g=\mathrm{diag}(a,1)$ 即可验）。由 $\operatorname{Im}(g(i))=\det(g)/|ci+d|^2$ 知正确形式为 $|\det(g)|^{1/2}|ci+d|^{-1}=|\operatorname{Im}(g(i))|^{1/2}$，译文按此。}但这确实成立，因为两边都在下列作用下不变：
\begin{itemize}
\item 从右的 $K_\infty=\SO_2(\R)$ 作用，
\item 从左的子群 $\begin{pmatrix}1&*\\0&1\end{pmatrix}\subseteq\GL_2(\R)$ 作用，
\item 从左的 $\GL_2(\R)$ 中对角矩阵的作用。
\end{itemize}
于是
\[
\int_{\Gamma A_{\GL_2}\backslash\GL_2(\R)}|\widetilde\phi_f(g)|^2\,\mathrm{d}g
=\int_{\Gamma\backslash\HH^\pm}|f(z)|^2|y|^k\frac{1}{y^2}\,\mathrm{d}x\wedge\mathrm{d}y
\]
（回忆我们是在 $\HH^\pm$ 的「那个」不变测度上积分）。由于 $\Gamma\backslash\HH^\pm$ 只有有限多个尖点，这引导到尖点处的局部陈述。通过共轭可以归约到考察 $\infty\in\Pone(\Q)$。把 $f$ 在 $\infty$ 处的 Fourier 展开写成
\[
f(z)=\sum_{n\in\Z}a_nq^{n/h},
\]
其中 $h\ge0$（依赖于 $\Gamma$）。则
\[
|y|^kf(z)=2\pi|\log(|q|)|^k\sum_{n\in\Z}a_nq^{n/h},
\]
因为 $y=\log(|q|)$。现在对 $0<r\le1$，
\[
\int_{0<|q|\le r}|\log(|q|)|^{k-2}\Bigl|\sum_{n\in\Z}a_nq^{n/h}\Bigr|^2\frac{i}{2}\,\mathrm{d}q\wedge\mathrm{d}\bar q<\infty,
\]
只要 $n\le0$ 时 $a_n=0$，即只要 $f$ 是尖点的\znote{被积函数在 $q\to0$ 时的行为由常数项 $a_0$ 决定：若 $a_0=0$，则 $\bigl|\sum_{n\ge1}a_nq^{n/h}\bigr|^2=O(|q|^{2/h})$，而 $\int_0^r|\log|q||^{k-2}|q|^{2/h}\,\mathrm{d}|q|$ 收敛；故这里真正用到的是常数项消失，仅「在尖点处全纯」不足以保证平方可积。[来源：原文第 22 页与本讲的整段推导]}。

因此，作为前面讨论的总结，映射
\[
M_k\subseteq H^0(\GL_2(\Q)\backslash(\GL_2(\Af)\times\HH^\pm),\omega^{\otimes k})\to C^\infty([\GL_2]),\qquad f\mapsto\widetilde\phi_f
\]
诱导典范包含
\[
\widetilde\Phi:S_k\hookrightarrow L^2([\GL_2]),
\]
% ===== PDF p.023 =====
其中 $S_k:=\varinjlim_KS_k(K)$。

我们可以作如下观察：
\begin{itemize}
\item $\GL_2(\Af)$ 作用在
\[
\begin{gathered}
H^0(\GL_2(\Q)\backslash(\GL_2(\Af)\times\HH^\pm),\omega^{\otimes k}),\quad
C^\infty(\GL_2(\Q)\backslash\GL_2(\A)),\\
C^\infty([\GL_2]),\quad L^2([\GL_2]),\quad M_k,\quad S_k
\end{gathered}
\]
上。
\item 态射
\[
\Phi:M_k\to C^\infty(\GL_2(\A)),\qquad f\mapsto\phi_f
\]
是 $\GL_2(\A)$-等变的，但
\[
M_k\to C^\infty([\GL_2]),\qquad f\mapsto\widetilde\phi_f=|\det|^{k/2}_{\mathrm{ad\acute{e}lique}}\phi_f
\]
不是（原因在于因子 $|\det|^{k/2}_{\mathrm{ad\acute{e}lique}}$）。
\item 更精确地，对 $g\in\GL_2(\Af)$ 与变动的 $K\subseteq\GL_2(\Af)$，同构
\[
M_k(K)\xrightarrow{\ g\ }M_k(g^{-1}Kg)\subseteq M_k
\]
通过 $M_k=\varinjlim_KM_k(K)$ 定义了作用 $M_k\xrightarrow{\ g\ }M_k$。
\item 对任意固定的 $K\subseteq\GL_2(\Af)$，$\GL_2(\Af)$ 不作用在 $M_k(K)$、$S_k(K)$ 等上\znote{这正是必须取归纳极限的原因：$\GL_2(\Af)$ 的作用只把 $K$ 送到共轭子群 $g^{-1}Kg$，因此单个 $M_k(K)$ 一般不稳定；只有 $\varinjlim_KM_k(K)$ 才承载 $G(\Af)$ 的作用。也正因如此，$M_k$ 不是某一个同余子群的模形式空间，而是「所有级」的模形式。[来源：原文第 22、23 页与本讲的构造动机]}。
\end{itemize}

$\GL_2(\Af)$-作用的存在，是我们从 $\Gamma$-模形式转向 $K$-模形式的主要动机。

于是我们看到接下来的任务：
\begin{itemize}
\item 把 $S_k$ 描述为 $\GL_2(\Af)$-表示。
\item 刻画 $S_k$ 在 $L^2([\GL_2])$ 中的像（或 $M_k$ 在 $C^\infty(\GL_2(\Q)\backslash\GL_2(\A))$ 中的像）。
\end{itemize}

希望模形式与自守表示（在 $L^2$ 意义下）的联系现在更清楚了。对 $f\in S_k$，可以取由 $f$ 的 $\GL_2(\A)$-平移生成的子空间的闭包，并希望它定义一个不可约酉表示。注意这样得到的表示自动落在 $L^2([\GL_2])$ 的离散部分 $L^2_{\mathrm{disc}}([\GL_2])$ 中。
